% Publication gate: retain the 87-matrix Lean claim only after the root
% worker obtains the complete PASS receipt for all 87 FourMatrix239 modules.
\subsection{Recovery after four missing neighbour counts}
\label{sec:four-erasure239}
The three-count result extends to four missing observations. The next
lemma supplies the additional arithmetic. Its hypothesis uses the
positions of the missing counts, independently of any snake ordering.

\begin{lemma}\label{lem:four-matrix239}
Let $p_0,p_1,p_2,p_3$ be distinct vertices of one parity in $Q_9$.
Let $M$ be the symmetric matrix with diagonal entries 128 and
\[
M_{ij}=\begin{cases}
-16,&\Ham(p_i,p_j)=2,\\
8,&\Ham(p_i,p_j)=4,\\
-10,&\Ham(p_i,p_j)=6,\\
35,&\Ham(p_i,p_j)=8
\end{cases}\qquad(i\ne j).
\]
If $w\in\mathbb Z^4$, $s\in\{-1,0,1\}^4$, and $Mw=315s$, then $w=0$.
\end{lemma}
\begin{proof}
We repeatedly use the following elementary divisibility fact:
if $Dx=315N$ with $D>0$ and integers $x,N$, then
$D/\gcd(D,315)$ divides $N$. If $|N|<D/\gcd(D,315)$, both $N$ and $x$
vanish. This also applies to each coordinate of an adjugate identity.

Join two missing positions when their distance is eight, so the
corresponding coefficient is 35. Translate one position to zero.
Its distance-eight neighbours are $J\oplus e_i$, where $J$ is the
all-ones word and $e_i$ has just its $i$th coordinate equal to one.
Two distinct such neighbours have distance two. The resulting graph
is therefore triangle-free. On four vertices its seven possible
unlabelled types are the empty graph, one edge, two disjoint edges,
a two-edge path with an isolated vertex, a three-leaf star, a
three-edge path, and a four-cycle. We treat these cases below.

\smallskip\noindent\emph{Empty graph.}
All coefficients are even. Parity forces $s=0$. If $m$ is the largest
absolute unknown, its row gives $128m\le48m$, hence $w=0$.

\smallskip\noindent\emph{One edge.}
Call its unknowns $x,y$ and the other two $z,t$. The last two right-hand
sides vanish by parity. Put $m=\max(|x|,|y|)$ and
$h=\max(|z|,|t|)$. Since every other coefficient has absolute value at
most 16, the equations imply
\[
112h\le32m,\qquad93m\le315+32h.
\]
Thus $587m\le2205$, giving $m\le3$ and $h\le6/7$. Integrality forces
$z=t=0$. The remaining coefficient-35 two-variable system vanishes by
the argument in Section~\ref{sec:three-erasure238}.

\smallskip\noindent\emph{Three-leaf star.}
Every leaf pair has coefficient $-16$. Subtracting two leaf rows gives
$144(y-z)=315(t-u)$. Since $|t-u|\le2$ and 144 divides neither 315 nor
630, the leaf unknowns and right-hand sides agree. The remaining
system has matrix
\[
\begin{pmatrix}128&105\\35&96\end{pmatrix}.
\]
Its determinant is 8613, with gcd 9 with 315. Hence 957 divides
$128t-35s$, whose absolute value is at most 163. It vanishes; since
$|s|,|t|\le1$, this forces $s=t=0$ and all unknowns zero.

\smallskip\noindent\emph{Three-edge path.}
Successive positions can be written
$0,J\oplus e_i,e_i\oplus e_j,J\oplus e_i\oplus e_j\oplus e_k$.
Distinctness and the absence of a closing distance-eight edge force
$i,j,k$ to be distinct. Thus the nonconsecutive distances are $2,2,6$,
and, in path order,
\[
M=\begin{pmatrix}
128&35&-16&-10\\35&128&35&-16\\
-16&35&128&35\\-10&-16&35&128
\end{pmatrix}.
\]
Adding and subtracting reversed pairs of rows yields the two matrices
\[
\begin{pmatrix}118&19\\19&163\end{pmatrix},\qquad
\begin{pmatrix}138&51\\51&93\end{pmatrix}.
\]
Their right-hand sides are 315 times integers of absolute value at
most two. Their determinants are 18873 and 10233, each with gcd 9 with
315. The reduced divisors 2097 and 1137 exceed the corresponding
adjugate numerator bounds 364 and 378. All sums and differences of
unknowns therefore vanish.

\smallskip\noindent\emph{Four-cycle.}
Opposite vertices have distance two. Subtracting their rows, again
with coefficient 144, makes opposite unknowns and right-hand sides
agree. The remaining matrix is
\[
\begin{pmatrix}112&70\\70&112\end{pmatrix}.
\]
Adding gives $182(x+y)=315(s+t)$. Since $\gcd(182,315)=7$ and
$|s+t|\le2$, we have $x+y=s+t=0$. The first row reduces to
$2x=15s$. Parity and $|s|\le1$ force $s=x=0$.

\smallskip\noindent\emph{Two-edge path and an isolated vertex.}
Put the centre at zero and the leaves at $J\oplus e_i,J\oplus e_j$.
The isolated position $q$ has weight $k\in\{2,4,6\}$. Its distance to
a leaf is $10-k$ when the corresponding coordinate belongs to $q$,
and $8-k$ otherwise. Excluding any additional distance-eight edge
gives the seven coupling triples $(\alpha,\beta,\gamma)$ in the table
below, up to interchange of the leaves.

The centre-and-leaves block $K$ and its adjugate are
\[
K=\begin{pmatrix}128&35&35\\35&128&-16\\35&-16&128\end{pmatrix},
\qquad
\operatorname{adj}K=
\begin{pmatrix}
16128&-5040&-5040\\-5040&15159&3273\\-5040&3273&15159
\end{pmatrix},
\]
with $\det K=1711584$. Put
$v=(\alpha,\beta,\gamma)^{\mathsf T}$ and $L=(\operatorname{adj}K)v$.
The isolated right-hand side is zero by parity. If $h$ is its unknown,
block elimination gives
\[
D h=-315L^{\mathsf T}(s,t,u),\qquad
D=128\det K-v^{\mathsf T}L.
\]
The seven exact calculations are
\begin{center}
\begin{tabular}{rrrr}
\toprule
$(\alpha,\beta,\gamma)$ & $D$ & $\gcd(D,315)$ & $\sum_i|L_i|$\\
\midrule
$(-16,-10,-10)$ & 214493184 & 9 & 364608\\
$(8,8,8)$ & 216981504 & 9 & 262656\\
$(8,-10,8)$ & 215926884 & 9 & 353052\\
$(8,-10,-10)$ & 212751360 & 45 & 679104\\
$(-10,-16,-16)$ & 211258368 & 9 & 489024\\
$(-10,8,-16)$ & 214263360 & 45 & 406224\\
$(-10,8,8)$ & 213497856 & 9 & 637632\\
\bottomrule
\end{tabular}
\end{center}
In every row, $D/\gcd(D,315)$ exceeds the last column. The scalar
product and $h$ vanish. The remaining three unknowns vanish by
Section~\ref{sec:three-erasure238}.

\smallskip\noindent\emph{Two disjoint edges.}
Write the first edge as $0,J\oplus e_i$ and the second as
$q,q\oplus J\oplus e_j$. Set $k=\operatorname{wt}(q)$, and let
$a,b\in\{0,1\}$ indicate membership of $i,j$ in the support of $q$.
For $i\ne j$, the four cross distances are
\[
\begin{pmatrix}k&8-k+2b\\8-k+2a&k+2-2(a+b)\end{pmatrix}.
\]
For $i=j$, the bottom-right entry is $k$ and $a=b$. Requiring all four
entries to belong to $\{2,4,6\}$, and interchanging the two edges or
their endpoints, leaves five symmetric cross-coefficient matrices:
\[
C_1=\begin{pmatrix}-16&-10\\-10&-16\end{pmatrix},\quad
C_2=\begin{pmatrix}8&8\\8&8\end{pmatrix},\quad
C_3=\begin{pmatrix}8&-10\\-10&8\end{pmatrix},
\]
\[
C_4=\begin{pmatrix}-16&-10\\-10&8\end{pmatrix},\qquad
C_5=\begin{pmatrix}8&8\\8&-10\end{pmatrix}.
\]
The cases $i=j$ give $C_1,C_2,C_3$; the cases $i\ne j$ give $C_4,C_5$
and endpoint permutations. For instance,
$\left(\begin{smallmatrix}-10&-16\\8&-10\end{smallmatrix}\right)$
becomes $C_4$ by interchanging the second edge's endpoints.

Put $B=\left(\begin{smallmatrix}128&35\\35&128\end{smallmatrix}\right)$.
The full matrix is $\left(\begin{smallmatrix}B&C\\C&B\end{smallmatrix}\right)$.
Adding its block equations gives $(B+C)X=315S$, where $X$ contains
sums of corresponding unknowns and $|S_i|\le2$. In the next table,
$q_\pm=\det(B\pm C)/\gcd(\det(B\pm C),315)$;
$R_+$ bounds each adjugate numerator for right-hand-side entries at
most two in absolute value, and $R_-$ does so for ternary entries.
\begin{center}
\begin{tabular}{rrrrr}
\toprule
$C$ & $q_+$ & $R_+$ & $q_-$ & $R_-$\\
\midrule
$C_1$ & 3973 & 274 & 297 & 189\\
$C_2$ & 5549 & 358 & 217 & 147\\
$C_3$ & 851 & 322 & 275 & 165\\
$C_4$ & 1623 & 322 & 339 & 189\\
$C_5$ & 4733 & 358 & 1759 & 165\\
\bottomrule
\end{tabular}
\end{center}
Since $q_+>R_+$, the sums of unknowns vanish, as do their right-hand
sides. The second block is now the negative of the first, so the first
block equation becomes $(B-C)(x,y)^{\mathsf T}=315(s,t)^{\mathsf T}$
with ternary $s,t$. Since $q_->R_-$, these unknowns also vanish.
This completes all seven cases.
\end{proof}

\begin{proposition}\label{prop:four-erasure239}
The neighbour counts at any 508 of the 512 vertices of $Q_9$ determine
a binary vertex selection uniquely. If the selected parity side is
known, any 252 of the 256 opposite-side counts suffice.
\end{proposition}
\begin{proof}
Let $f$ be the difference of two binary selection indicators, and
write $Af(v)=\sum_{u\sim v}f(u)$ for its neighbour sums. Then $f$ is
ternary and $Af$ vanishes outside the four missing positions.
If those positions have mixed parity, restrict $f$ to each parity
side. Adjacency sends each restriction to the opposite side, where
its neighbour sums have support at most three. The three-erasure
result makes both restrictions zero.

If the four positions have one parity, apply the inverse formula at
their antipodes. Its diagonal coefficient is 128; at pairwise distances
$2,4,6,8$ its off-diagonal coefficients are $-16,8,-10,35$.
Lemma~\ref{lem:four-matrix239} makes all four exceptional neighbour
sums zero. The inverse then gives $f=0$ everywhere. For the parity
specialization, all 256 counts on the selected side are already zero,
so knowing 252 opposite-side counts leaves just four observations
unspecified.
\end{proof}

\begin{proposition}\label{prop:eight-erasure239}
Let $E\subseteq V(Q_9)$ have at most eight vertices, with no two at
Hamming distance eight. The neighbour counts outside $E$ determine a
binary selection uniquely. The number eight is sharp under this
condition on $E$.
\end{proposition}
\begin{proof}
Subtract two binary selection indicators and call their difference
$f$. Suppose $Af$ vanishes outside $E=\{p_1,\ldots,p_k\}$.
The case $k=0$ follows from the inverse. Otherwise put $w_i=Af(p_i)$
and apply the inverse at the antipodes of these positions. Its diagonal
entries are 128. Excluding distance eight leaves off-diagonal entries
in $\{0,-10,8,-16\}$, all even and of absolute value at most 16.
Every equation has right-hand side $315s_i$ with
$s_i\in\{-1,0,1\}$. Parity therefore forces every $s_i=0$.
If $M=\max_i|w_i|$, a row attaining the maximum gives
\[
128M\le16(k-1)M\le112M.
\]
Thus every exceptional count is zero, and the inverse gives $f=0$.
This proof permits both parities among the missing positions.

For sharpness, let $E$ consist of the nine neighbours of a single
vertex $v$. Every pair in $E$ has distance two. The empty selection
and the selection $\{v\}$ have identical counts outside $E$.
Equivalently, the nine-position antipodal matrix is $144I-16J$,
where $J$ is the all-ones matrix, and it annihilates the all-ones vector.
\end{proof}

Proposition~\ref{prop:eight-erasure239} is a written consequence of the
cube inverse, parity and the displayed scalar inequality. Its geometric
hypothesis is essential to the argument: for arbitrary missing
positions, Proposition~\ref{prop:four-erasure239} gives the established
four-count result. Any nonzero difference hidden at five through eight
positions must include a distance-eight pair among those positions.

The finite audit independently covers the geometry. After translating
one position to zero, record the multiplicities of the eight possible
three-bit coordinate columns. The $\binom{16}{7}=11440$ compositions of
nine columns contain 1149 valid ordered descriptions, giving 105
four-point set orbits under cube isometries and 87 distinct distance
matrices up to point relabelling. Exact adjugate calculations over all
$87\cdot81=7047$ ternary right-hand sides find no nonzero integer
solution, with no bound on the unknowns. A separate direct enumeration
of all $\binom{255}{3}=2731135$ triples of nonzero even-weight vertices
recovers exactly the same 87 matrix classes~\cite{research239}.

All 87 literal matrix implications have passed Lean checking. Each
checks its adjugate identities and finite divisibility calculation,
then concludes that every integer solution is zero. The geometric
classification, its connection to the matrix catalogue, and the
mixed-parity reduction in Proposition~\ref{prop:four-erasure239} remain
written arguments supported by the independent native coverage audit.
The complete four-erasure statement on the cube has no single assembled
Lean theorem in this record.

Two distinct binary selections must consequently differ in their
neighbour counts at at least five positions. A one-vertex change alters
nine counts, leaving the minimum number between five and nine. These
recovery results impose no connectedness condition and yield no longer
snake or new global snake-length bound.
